\documentclass[../../../uem_mat_mei_q.tex]{subfiles}

\begin{document}
    \ref{q:uem_mat_mei_2025_icom_04_q:1}\,～\,\ref{q:uem_mat_mei_2025_icom_04_q:3}\,%
    については所定の解答欄（表面）に解答をマークせよ．%
    \ref{q:uem_mat_mei_2025_icom_04_q:4}\,%
    については所定の解答欄（裏面）に解答経過と答をともに記せ．

    問題文中の$\myboxb{ノ}$は解答が1ケタの数であることを，\\
    $\myboxb{ヌメ}$などは解答が2ケタの数であることを表している．

    なお，分数形で解答する場合，それ以上約分できない形で答えること．

    \vspace{7pt}
    下の図のような，東西に延びる5本の道と，南北に延びる5本の道が交差して，%
    碁盤目のように区分された街がある．なお，2本の道が交わるところをポイントと呼ぶことにする．%
    図中に示すように，街の中にはポイントA，B，C，D，Eがあり，%
    ポイントAは街の南西角に，ポイントBは北東角にある．

    この街でXさんとYさんが，次のルールに従って移動するものとする．

    \begin{itemize}[label={\textbullet}, leftmargin=2\zw, labelsep=.4\zw, topsep=7pt, itemsep=1pt]
        \item Xさんは東または北にしか進まない．
        \item Yさんは西または南にしか進まない．
        \item XさんもYさんも，各ポイントで同時にそれぞれサイコロを振り，%
            出た目に応じて下の表に示す方向に1つ隣りのポイントまで進む．%
            ただし，一方向にしか進めないときは，サイコロの目にかかわらず，確率1でその方向に進む．
    \end{itemize}
    \begin{minipage}{\linewidth}
        \centering
            \includegraphics%
                {\subfix{fig_uem_mat_mei_2025_icom_04/fig_uem_mat_mei_2025_icom_04_q_01.pdf}}\\
            \vspace{7pt}
            \includegraphics%
                {\subfix{fig_uem_mat_mei_2025_icom_04/fig_uem_mat_mei_2025_icom_04_q_02.pdf}}%
    \end{minipage}
    
    \vspace{9pt}
    このとき，以下の問に答えよ．
    \newpage
    \begin{enumerate}
        \item%
            XさんがポイントAにいるとする．このとき，XさんがポイントC，D，Eのいずれも%
            通らずにポイントBに到着する経路はすべてで$\myboxb{ヌネ}$通りある．
            \label{q:uem_mat_mei_2025_icom_04_q:1}
        \item%
            XさんがポイントAにいるとする．このとき，Xさんが，ポイントCを通り，%
            かつポイントDを通らず，さらにポイントEを通ってポイントBに到着する確率は%
            $\myfracc{\myboxb{ノ}}{\myboxb{ハヒ}}$である．
        \item%
            XさんがポイントCに，YさんがポイントEにいるとする．この後，%
            XさんとYさんが出会う確率は$\myfracc{\myboxb{フヘ}}{\myboxb{ホマ}}$である．
            \label{q:uem_mat_mei_2025_icom_04_q:3}
        \item%
            XさんがポイントAに，YさんがポイントBにいるとする．この後，%
            XさんとYさんがポイントD以外で出会う確率を求めよ．
            \label{q:uem_mat_mei_2025_icom_04_q:4}
    \end{enumerate}

\end{document}